\subsection{THE MULTIPLICATIVE GROUP U}

THEOREM 1. The (multiplicative) topological group U of complex numbers of absolute value 1 is isomorphic to the (additive) topological group T of real numbers mod 1.

\(\mathbf{U} = \mathbf{S}_{1}\) is compact and connected, and has a neighbourhood of the identity element \(+1\) homeomorphic to an open interval of \(\mathbf{R}\) (Chapter VI, § 2, no. 4, Proposition 5); the theorem is therefore a consequence of the topological characterization of \(\mathbf{T}\) given in Chapter V, § 3 (Theorem 2).

COROLLARY. The multiplicative group \(\mathbf{C}^*\) of non-zero complex numbers is isomorphic to the group \(\mathbf{R} \times \mathbf{T}\) (cf.~§ 1. no. 3, Proposition 1).

Remark. The isomorphism of the groups \(C^{*}\) and \(R \times T\) implies the existence of roots of every ``binomial equation'' \(z^{n} = a\) in the field C. Using this fact and the local compactness of C we can obtain another proof of the theorem of d'Alembert-Gauss (Exercise 2).

There are only two distinct isomorphisms of the group \(\mathbf{T}\) onto the group \(\mathbf{U}\) ; for if \(g\) , \(g'\) are two isomorphisms of \(\mathbf{T}\) onto \(\mathbf{U}\) , and if \(h'\) is the inverse of \(g'\) , then \(h' \circ g\) is an automorphism of \(\mathbf{T}\) , and therefore (Chapter VII, § 2, no. 4, Proposition 6) we have identically either \(g'(x) = g(x)\) or \(g'(x) = g(-x)\) . We may always assume that \(g\) is such that \(i\) is the image under \(g\) of the class mod \(1\) of the point \(\frac{1}{4}\) ; then, if \(\varphi\) denotes the canonical homomorphism of \(\mathbf{R}\) onto \(\mathbf{T}\) , every strict morphism of the additive group \(\mathbf{R}\) onto the multiplicative group \(\mathbf{U}\) is of the form \(x \to g(\varphi(x/a))\) , where \(a\) is a real number \(\neq o\) (Chapter VII, § 2, no. 3, Proposition 4); note that the interval \([- \frac{1}{2} |a|, \frac{1}{2} |a|]\) is the largest symmetric open interval of \(\mathbf{R}\) which is mapped one-to-one onto its image by this strict morphism, and that we have \(g(\varphi(\frac{1}{4})) = i\) . We shall denote the homomorphism \(x \to g(\varphi(x))\) by \(x \to e(x)\) ; every strict morphism of \(\mathbf{R}\) onto \(\mathbf{U}\) is therefore of the form \(x \to e(x/a)\) , where \(a \neq 0\) . The function \(e(x)\) is continuous on \(\mathbf{R}\) , complex valued, and satisfies the identities

\begin{enumerate}
\def\labelenumi{(\arabic{enumi})}
\tightlist
\item
\end{enumerate}

\[
\left| e (x) \right| = 1,\tag{2}
\]

\[
\boldsymbol {e} (x + y) = \boldsymbol {e} (x) \boldsymbol {e} (y),
\]

together with the relations

\[
e (0) = \mathrm{I}, \quad e (\frac {1}{4}) = i, \quad e (\frac {1}{2}) = - \mathrm{I}, \quad e (\frac {3}{4}) = - i, \quad e (\mathrm{I}) = \mathrm{I}. \tag {3}
\]

From (1) and (2) it follows that

\[
e (- x) = \frac {1}{e (x)} = \overline {{{{e (x)}}}},\tag{4}
\]

and from (2) and (3) that

\[
\begin{array}{l} e (x + \frac {1}{4}) = i e (x), \quad e (x + \frac {1}{2}) = - e (x), \\ e (x + \frac {3}{4}) = - i e (x), \quad e (x + 1) = e (x). \end{array}
\]

Thus the function \(e(x)\) is periodic and has 1 as a principal period.

Remark. The mapping \(x + iy \to e^x e(y)\) is a strict morphism of the additive group \(\mathbf{C}\) onto the multiplicative group \(\mathbf{C}^*\) , and its restriction to a suitable neighbourhood of \(\mathbf{o}\) is a local isomorphism of \(\mathbf{C}\) with \(\mathbf{C}^*\) . Consequently (Chapter VII, § 2, no. 3) every strict morphism of \(\mathbf{C}\) onto \(\mathbf{C}^*\) is of the form \(x + iy \to e^{\alpha x + \beta y} e(\gamma x + \delta y)\) , where \(\alpha, \beta, \gamma, \delta\) are any real numbers such that \(\alpha \delta - \beta \gamma \neq \mathbf{o}\) . We shall see later that there is just one of these homomorphisms, denoted by \(z \rightarrow e^{z}\) , such that

\[
\lim _ {z \rightarrow 0} \frac {e ^ {z} - 1}{z} = 1;
\]

and the restriction of this homomorphism to the real axis is the same as \(e^{x}\) (whence the notation).

